Using Kepler's Third law and the semi-major axis ($a = 30.070$\,au) of the
orbit from the table of constants, the sidereal period of Neptune's orbit is:
\[
  P_{\mathrm{N}} = (a^3)^{\frac{1}{2}} = \sqrt{27189.4}
    = 164.89\ \mathrm{years} \hfill
\]
\hfill(1 point)

The rest of the calculation can be done in (at least) two ways:

\medskip
\noindent\textbf{(1) `day counting'/`date drift' solution:}

Since Neptune is an outer planet (relative to Earth), the synodic period $S$ in
years is given by:
\[
  \frac{1}{S} = \frac{1}{1\ \mathrm{year}} - \frac{1}{P_{\mathrm{N}}}
\]
\hfill(1 point)
\begin{align*}
  &= 1 - \frac{1}{164.89} = 1 - 0.006065 = 0.99394\\
  &\Longrightarrow S = 1/0.99394 = 1.006102\ \mathrm{years}
\end{align*}
\hfill(1 point)
\[
  \Longrightarrow 1.006102 \times 365.2422 = 367.4708\ \mathrm{solar\ days}
\]

Therefore the date of opposition drifts by:
\[
  D = 367.4708 - 365.2422 = 2.2286\ \mathrm{days/year}
\]
\hfill(1 point)

The approximate date of the northern spring equinox is 20 March. The number of
days between 21 September and 20 March $= 185$ days, therefore year of desired
opposition is:
\[
  2024 - (185/D) = 2024 - 83 = 1941
\]
\hfill(1 point)

If the student takes 21 March as the spring equinox, they get
$184/D = 82.5\,\mathrm{yr} \Longrightarrow 1942$.

\medskip
\noindent\textbf{(2) `angular drift' solution (more accurate)}

As before, the synodic period of Neptune can be derived as:
\[
  \frac{1}{S} = \frac{1}{1\ \mathrm{year}} - \frac{1}{P_{\mathrm{N}}}
\]
\hfill(1 point)
\[
  = 1 - \frac{1}{164.89} = \frac{163.89}{164.89}
  \Longrightarrow S = \frac{164.89}{163.89}\ \mathrm{years}
\]
Therefore the ecliptic longitude of Neptune at opposition drifts by
$(360^\circ/163.89)/\mathrm{year}$. We want to find how many years it takes for
the longitude of opposition to move by $180^\circ$, i.e.\ for what $t$:
\[
  t \times \frac{360^\circ}{163.89} = 180^\circ.
\]
\hfill(2 points)
\[
  \Longrightarrow t = 163.89/2 = 81.95\ \mathrm{years}
  \Longrightarrow 2024 - 82 = 1942
\]
\hfill(1 point)

We accept 1941 or 1942 for full points for calculations using the assumptions
in the question. If the student uses some other method which is conceptually
correct and results in 1943 (the true answer) they should also get full
points.

\medskip
Table of spring equinoxes and oppositions of Neptune:

\begin{center}
\begin{tabular}{|c|c|c|c|r|}
\hline
Year & Equinox (UT) & Opposition (UT) & Coordinates of Neptune
  & $\Delta T$ [days]\\
\hline
1940 & Mar 20 18:42 & Mar 14 21:08
  & $11^{\mathrm{h}}40^{\mathrm{m}}$ $+3^\circ 28'$ & $5.9$\\
1941 & Mar 21 00:20 & Mar 17 07:40
  & $11^{\mathrm{h}}49^{\mathrm{m}}$ $+2^\circ 39'$ & $3.7$\\
1942 & Mar 21 06:11 & Mar 19 18:12
  & $11^{\mathrm{h}}57^{\mathrm{m}}$ $+1^\circ 50'$ & $1.5$\\
1943 & Mar 21 12:03 & Mar 22 04:51
  & $12^{\mathrm{h}}04^{\mathrm{m}}$ $+1^\circ 00'$ & $-0.7$\\
1944 & Mar 20 17:49 & Mar 23 15:29
  & $12^{\mathrm{h}}13^{\mathrm{m}}$ $+0^\circ 11'$ & $-2.9$\\
\hline
2024 & Mar 20 03:06 & Sep 21 00:16
  & $23^{\mathrm{h}}55^{\mathrm{m}}$ $+1^\circ 56'$ & $-184.9$\\
\hline
\end{tabular}
\end{center}

Taking into account all effects, the last opposition closest to the spring
equinox was actually in 1943. 1942 results from the assumptions made in the
question.
